% Options for packages loaded elsewhere
\PassOptionsToPackage{unicode}{hyperref}
\PassOptionsToPackage{hyphens}{url}
%
\documentclass[
  11pt,
]{article}
\usepackage{amsmath,amssymb}
\usepackage{iftex}
\ifPDFTeX
  \usepackage[T1]{fontenc}
  \usepackage[utf8]{inputenc}
  \usepackage{textcomp} % provide euro and other symbols
\else % if luatex or xetex
  \usepackage{unicode-math} % this also loads fontspec
  \defaultfontfeatures{Scale=MatchLowercase}
  \defaultfontfeatures[\rmfamily]{Ligatures=TeX,Scale=1}
\fi
\usepackage{lmodern}
\ifPDFTeX\else
  % xetex/luatex font selection
\fi
% Use upquote if available, for straight quotes in verbatim environments
\IfFileExists{upquote.sty}{\usepackage{upquote}}{}
\IfFileExists{microtype.sty}{% use microtype if available
  \usepackage[]{microtype}
  \UseMicrotypeSet[protrusion]{basicmath} % disable protrusion for tt fonts
}{}
\makeatletter
\@ifundefined{KOMAClassName}{% if non-KOMA class
  \IfFileExists{parskip.sty}{%
    \usepackage{parskip}
  }{% else
    \setlength{\parindent}{0pt}
    \setlength{\parskip}{6pt plus 2pt minus 1pt}}
}{% if KOMA class
  \KOMAoptions{parskip=half}}
\makeatother
\usepackage{xcolor}
\usepackage[margin=25mm]{geometry}
\setlength{\emergencystretch}{3em} % prevent overfull lines
\providecommand{\tightlist}{%
  \setlength{\itemsep}{0pt}\setlength{\parskip}{0pt}}
\setcounter{secnumdepth}{-\maxdimen} % remove section numbering
\usepackage{amssymb,amsmath}
\ifLuaTeX
  \usepackage{selnolig}  % disable illegal ligatures
\fi
\usepackage{bookmark}
\IfFileExists{xurl.sty}{\usepackage{xurl}}{} % add URL line breaks if available
\urlstyle{same}
\hypersetup{
  pdftitle={A sharp Gaussian first-moment inequality for every prescribed mass vector},
  pdfauthor={Yongxian Zhang},
  hidelinks,
  pdfcreator={LaTeX via pandoc}}

\title{A sharp Gaussian first-moment inequality for every prescribed
mass vector}
\author{Yongxian Zhang\thanks{School of Computer Science and Engineering,
South China University of Technology, Guangzhou, China.
Email: \href{mailto:mxymmxym1@gmail.com}{mxymmxym1@gmail.com}.
ORCID: \href{https://orcid.org/0009-0000-3864-3536}{0009-0000-3864-3536}.}}
\date{8 October 2026}

\begin{document}
\maketitle

\section{Abstract}\label{abstract}

We establish a sharp first-moment inequality for every measurable
simplex-valued function on Gaussian space with prescribed positive mean.
The label metric is the negative Hessian of the Gaussian multi-bubble
perimeter profile. The equality functions are exactly the indicators of
the translated regular simplex partition with those probabilities,
extended cylindrically. The result holds for every number of labels and
every ambient dimension. Its proof derives a global tangent bound for
the squared optimal Gaussian assignment value from the established
multi-bubble perimeter theorem and a radial differential inequality. An
exact nonnegative integral measures its deficit. At equal masses this
recovers the sharp ordinary squared-centroid sum; at unequal masses the
metric is generally anisotropic. We do not assert the false unequal-mass
Euclidean simplex conjecture or the full functional multi-bubble
isoperimetric inequality. The profile Hessian identity is prior work;
our statement is a bound for all functions, not only perimeter
minimizers.

\section{1. Model profile and the sharp functional
theorem}\label{model-profile-and-the-sharp-functional-theorem}

Let \(k\ge2\), \(n=k-1\), \(\mathbf1=(1,\ldots,1)^T\),
\(P=I-\mathbf1\mathbf1^T/k\), and \(Q_*=P/n\). Fix centered rows
\(v_i^*\) of a \(k\times n\) matrix \(V_*\) with \(V_*V_*^T=Q_*\). Their
pairwise distances are \(\ell_* =\sqrt{2/n}\).

Fix a probability vector \(p=(p_1,\ldots,p_k)\) with every \(p_i>0\).
There are unique prices \(\lambda_p\) in the gauge
\(\sum_i\lambda_{p,i}=0\) such that

\[
 \Omega_i(p)=\{x:v_i^*\cdot x-\lambda_{p,i}
                   \ge v_j^*\cdot x-\lambda_{p,j}\ \forall j\}
                                                               \tag{1}
\]

has standard Gaussian probability \(p_i\) in \(\mathbb R^n\). Their
existence, uniqueness and smooth dependence are proved in Section 2. The
score matrix maps \(\mathbb R^n\) onto \(\mathbf1^\perp\), so
\(\lambda_p=V_*w_p\) for a unique \(w_p\). Thus (1) is a central regular
simplex fan translated by \(w_p\).

Write \(A_{ij}^p\) for the ordinary Gaussian interface areas of this
model. Define two label Laplacians, annihilating \(\mathbf1\):

\[
 (H_pu)_i=\sum_{j\ne i}A_{ij}^p(u_i-u_j),\qquad
 L_p=H_p/\ell_*.
\]

All model facets have positive area, so both matrices are positive on
\(\mathbf1^\perp\). The notation \(H_p^+,L_p^+\) means inversion on this
subspace and zero on the constant line. Put

\[
 I(p)=\sum_{i<j}A_{ij}^p,\qquad
 c_p=\frac{\operatorname{tr}L_p}{n}=\ell_*I(p)>0.          \tag{2}
\]

By Milman--Neeman {[}MN22, Theorem 1.1{]}, \(I\) is the Gaussian
\(k\)-cell perimeter profile in dimension \(n=k-1\). Perimeter is
normalized as half the sum of individual cell perimeters, or
equivalently the sum of ordinary pairwise interfaces. Its Hessian on the
tangent space \(\mathbf1^\perp\) satisfies the established identity

\[
                         \nabla^2 I(p)=-H_p^+.           \tag{3}
\]

This identity is {[}MN22, Proposition 2.4{]}; Section 2 also derives its
normalization from the price calculus.

\textbf{Theorem 1.} Let \(d\ge1\) and let \(f_i:\mathbb R^d\to[0,1]\) be
measurable, with \(\sum_i f_i=1\) almost everywhere and
\(\int f_i\,d\gamma_d=p_i\). Let the \(k\times d\) moment matrix \(B\)
have rows \(b_i=\int xf_i(x)\,d\gamma_d(x)\). Then

\[
 \operatorname{tr}(B^TH_p^+B)\le I(p),                    \tag{4}
\]

or equivalently

\[
 -\sum_{a=1}^d B_{\cdot,a}^T\nabla^2 I(p)B_{\cdot,a}\le I(p).
                                                               \tag{5}
\]

If \(d\ge k-1\), equality holds exactly when the \(f_i\) are the
ordinary winning indicators (1) in an orthogonally embedded
\((k-1)\)-dimensional subspace, with its orthogonal complement
unrestricted, up to null sets and label symmetries preserving the fixed
mass vector. If \(d<k-1\), the inequality is strict for every such
function; no sharp fixed-lower-dimensional bound is claimed.

This is a universal functional first-moment bound. It applies to all
measurable partitions and fractional labels of mean \(p\), rather than
only to variations of perimeter minimizers. The profile identity (3) and
the perimeter theorem are explicitly imported prior results. The finite
literature check {[}LS{]} found no equivalent all-function statement in
the checked sources; it does not establish worldwide novelty.

\section{2. Prescribed-mass Gaussian score
calculus}\label{prescribed-mass-gaussian-score-calculus}

For every symmetric positive-semidefinite covariance \(Q\) with
\(Q\mathbf1=0\), define

\[
 \mathcal C_p(Q)=\inf_{\lambda\in\mathbb R^k}
       \left\{p\cdot\lambda+\mathbb E\max_i(Y_i-\lambda_i)\right\},
                 \qquad Y\sim N(0,Q).                    \tag{6}
\]

Adding a common price has no effect. In the gauge \(\min\lambda_i=0\),
the expression is at least \(p_{\min}\max\lambda_i\): a zero-priced
mean-zero score makes the expected maximum nonnegative. Hence the
infimum is attained. Pointwise comparison with the \(p\)-weighted
average of scores also shows \(\mathcal C_p(Q)\ge0\).

For distinct score rows \(v_i\), Gaussian pairwise ties are null. The
price gradient is \(p_i-\gamma_n(C_i)\), where
\(C_i=\{x:v_i\cdot x-\lambda_i\ge v_j\cdot x-\lambda_j\ \forall j\}\).
Minimizing prices therefore give precisely the masses \(p_i\). Their
winning assignment attains the score maximum pointwise. For any feasible
ordinary or fractional labels \(f\) with mean \(p\), and any centered
score matrix \(M\), the same pointwise comparison at every price gives

\[
               \langle M,B\rangle\le\mathcal C_p(MM^T).  \tag{7}
\]

No assignment interpretation at duplicated scores is needed.

On the entire closed covariance cone,

\[
 \mathcal C_p(sQ)=\sqrt{s}\,\mathcal C_p(Q)\quad(s\ge0),
                  \qquad\mathcal C_p\text{ is continuous}.       \tag{8}
\]

Scale scores and prices to prove homogeneity, and note
\(\mathcal C_p(0)=0\). For continuity couple \(Q^{1/2}G\) and
\(\widetilde Q^{1/2}G\). The difference of their maxima at any price is
at most \(|(Q^{1/2}-\widetilde Q^{1/2})G|\). Its expectation bounds the
two infima uniformly in prices. Continuity of the matrix square root
therefore covers singular endpoints as well.

Suppose \(Q\) is positive on \(\mathbf1^\perp\), and write \(Q=MM^T\)
with centered rows in \(\mathbb R^n\). They are affinely independent.
The \(n\) score differences against one label give nonsingular affine
coordinates. They realize a top tie between any chosen pair with all
remaining scores strictly lower. Consequently every pair facet has
positive Gaussian area. Define its length \(\ell_{ij}=|v_i-v_j|\), area
\(A_{ij}\), weight \(w_{ij}=A_{ij}/\ell_{ij}\), and weighted Laplacian
\(L\).

Gaussian integration by parts, with inward facet normal
\((v_i-v_j)/\ell_{ij}\), gives the winning-cell moments \(B(M)=LM\). The
additional boundary terms after truncation with balls tend to zero
exponentially. Since the winning masses are \(p\), price terms cancel,
so

\[
 \mathcal C_p(Q)=\operatorname{tr}(LQ)=\sum_{i<j}w_{ij}\ell_{ij}^2,
                 \qquad\operatorname{tr}L=2\sum_{i<j}w_{ij}.      \tag{9}
\]

\textbf{Lemma 2.} On this positive cone \(\mathcal C_p\) is smooth and

\[
                 d\mathcal C_p(Q)[D]=\tfrac12\operatorname{tr}(LD)
                                                               \tag{10}
\]

for every symmetric centered covariance direction \(D\).

\emph{Proof.} In the affine score-difference coordinates, winning
regions are fixed polyhedral cones. Gaussian density and
linear-integrand derivatives are locally dominated by a polynomial times
a Gaussian, so these integrals are smooth. The price Hessian is \(L\): a
price gap moves the pair facet by that gap divided by \(\ell_{ij}\). All
pair weights are positive, making \(L\) positive on the zero-sum gauge
at every price. Thus balancing prices are unique in that gauge, and the
implicit function theorem makes them smooth.

Write \(D=MAM^T\) with \(A\) symmetric; this representation exists
because \(M\) maps onto \(\mathbf1^\perp\). Realize its covariance path
by \(M(t)=M(I+tA)^{1/2}\). The envelope score derivative is \(B(M)\) and
\(M'(0)=MA/2\). Formula (9) then gives
\(d\mathcal C_p(Q)[D]=\operatorname{tr}(B(M)^TMA)/2
=\operatorname{tr}(LD)/2\). \(\square\)

This proves the assertions about model prices in Section 1. In the model
\(\ell_{ij}=\ell_*\); (9) gives
\(\mathcal C_p(Q_*)=\ell_*I(p)=\operatorname{tr}L_p/n=c_p\), including
all factors in (2). For completeness the derivative of the mass vector
with respect to model prices is \(-L_p\) on the tangent gauge. The
envelope derivative of \(c_p\) with respect to \(p\) is \(\lambda_p\).
Therefore \(\nabla^2 c_p=-L_p^+\), and
\(\nabla^2 I=-L_p^+/\ell_*=-H_p^+\), as in the prior identity (3).

\section{3. Perimeter and a global covariance tangent
bound}\label{perimeter-and-a-global-covariance-tangent-bound}

Let \(Q\) be positive on \(\mathbf1^\perp\) and of trace one. Its
balanced winning cells have masses \(p\) in dimension \(n=k-1\),
precisely the range \(k\le n+1\) of the imported multi-bubble theorem.
Write their perimeter \(S=\sum_{i<j}A_{ij}\). Thus \(S\ge I(p)\).
Weighted Cauchy--Schwarz and (9) give

\[
 S^2=\left(\sum_{i<j}w_{ij}\ell_{ij}\right)^2
 \le\left(\sum_{i<j}w_{ij}\right)
                    \left(\sum_{i<j}w_{ij}\ell_{ij}^2\right)
 =\tfrac12\mathcal C_p(Q)\operatorname{tr}L.
\]

Using \(I(p)=c_p\sqrt{n/2}\), we obtain

\[
                   \mathcal C_p(Q)\frac{\operatorname{tr}L}{n}
                              \ge c_p^2.                \tag{11}
\]

Every weight is positive, so equality in weighted Cauchy requires all
pair lengths equal. A centered trace-one score Gram with this property
is \(Q_*\): double-centering its squared-distance matrix gives
\(Q=(\ell^2/2)P\), and trace one gives \(\ell^2=2/n\).

\textbf{Theorem 3.} For every centered positive-semidefinite covariance
\(Q\),

\[
                  \mathcal C_p(Q)^2\le c_p\operatorname{tr}(L_pQ).
                                                               \tag{12}
\]

Equality holds exactly when \(Q=sQ_*\) for some \(s\ge0\).

\emph{Proof.} First normalize \(\operatorname{tr}Q=1\), and put
\(Q_t=Q_*+t(Q-Q_*)\) for \(0\le t<1\). This is positive on
\(\mathbf1^\perp\) even when \(Q\) is singular. Write
\(C(t)=\mathcal C_p(Q_t)\) and \(L_t\) for its facet matrix. Formulas
(9)--(10) imply

\[
              tC'(t)=\tfrac12\left(C(t)-\operatorname{tr}L_t/n\right).
                                                               \tag{13}
\]

Let \(h(t)=C(t)^2-c_p^2\) and
\(D_p(t)=C(t)\operatorname{tr}L_t/n-c_p^2\ge0\). Then exactly

\[
                     (h(t)/t)'=-D_p(t)/t^2\le0.           \tag{14}
\]

The initial derivative need not vanish. By (10) its value is

\[
 \lim_{t\downarrow0}h(t)/t=h'(0)
                      =c_p\operatorname{tr}(L_p(Q-Q_*)).  \tag{15}
\]

At one the quotient tends to \(\mathcal C_p(Q)^2-c_p^2\) by (8).
Integrating (14) on compact subintervals and taking both endpoint limits
therefore gives the exact identity

\[
 c_p\operatorname{tr}(L_pQ)-\mathcal C_p(Q)^2
                         =\int_0^1\frac{D_p(t)}{t^2}\,dt\ge0.
                                                               \tag{16}
\]

This nonnegative improper integral is finite because both quotient
limits are finite. Smoothness near zero also gives \(D_p(t)=O(t^2)\).
Here \(c_p\operatorname{tr}(L_pQ_*)=c_p^2\) was used to simplify the
endpoint difference. This proves (12) at trace one.

If equality holds, continuity and nonnegativity make \(D_p(t)=0\) for
every \(0<t<1\). The perimeter/Cauchy chain consists of equalities, and
all pair lengths are equal, forcing \(Q_t=Q_*\). Any such \(t>0\)
implies \(Q=Q_*\). Conversely the model has equality. For nonzero \(Q\)
of trace \(s>0\), apply the result to \(Q/s\) and use homogeneity (8).
At \(Q=0\) both sides vanish. This gives all stated equality rays,
including singular endpoints. \(\square\)

Writing \(W(t)=\sum_{i<j}w_{ij}(t)\), the integrand has the
decomposition

\[
 D_p(t)=\frac2n\left[
  W(t)\sum_{i<j}w_{ij}(t)\ell_{ij}(t)^2-S(t)^2
                                    +S(t)^2-I(p)^2\right].       \tag{17}
\]

Both terms are nonnegative: weighted edge-length variance and
squared-perimeter deficit. No metric stability rate or full covariance
concavity is asserted.

\section{4. Moment duality, sharpness and all equality
cases}\label{moment-duality-sharpness-and-all-equality-cases}

For the functions in Theorem 1 the rows of \(B\) sum to zero. Put
\(E=\operatorname{tr}(B^TL_p^+B)\ge0\) and \(M=L_p^+B\). The feasible
score inequality (7) and Theorem 3 imply

\[
 E=\langle M,B\rangle\le\mathcal C_p(MM^T),\qquad
 \mathcal C_p(MM^T)^2\le c_p\operatorname{tr}(L_pMM^T)=c_pE.       \tag{18}
\]

The last equality uses \(L_p^+L_pL_p^+=L_p^+\). Thus \(E^2\le c_pE\),
giving \(E\le c_p\). Since \(H_p=\ell_*L_p\) and \(c_p=\ell_*I(p)\),
this is exactly (4), and the tangent-space identity (3) gives (5).

If \(E=c_p>0\), all comparisons in (18) are equalities and
\(\mathcal C_p(MM^T)=c_p\). Theorem 3 gives \(MM^T=sQ_*\); homogeneity
then gives \(\sqrt{s}c_p=c_p\), so \(s=1\). Thus \(M\) is a trace-one
regular simplex score list. Equality in the feasible assignment means
its nonnegative pointwise maximum deficit has zero integral. Its scores
have null ties, so the fractional labels must be the unique winning
indicators. Uniqueness of the prescribed-mass prices identifies exactly
(1) after an orthogonal embedding. This requires \(d\ge n\); below that
dimension equality is impossible. If \(E=0\) the inequality is strict.

Conversely model flux gives \(B_*=L_pV_*\), whence

\[
 \operatorname{tr}(B_*^TL_p^+B_*)
         =\operatorname{tr}(V_*^TL_pV_*)
         =\operatorname{tr}(L_pQ_*)=c_p.
\]

Orthogonal and cylindrical extensions preserve moments and masses. This
proves Theorem 1, its sharpness and every equality assertion. The proof
includes \(k=2\), where the normalized covariance space is a singleton.
The omitted \(k=1\) case has one cell and zero moments. No statement at
a zero prescribed probability is needed.

\section{5. Equal masses and the scope of the
result}\label{equal-masses-and-the-scope-of-the-result}

At \(p_i=1/k\), symmetry gives \(L_p=c_kP\), where
\(c_k=\mathbb E\max_{i\le k}Z_i/\sqrt{k-1}\) for independent standard
normals. Theorem 1 then reduces to

\[
             \sum_i|b_i|^2\le
                   \frac{(\mathbb E\max_{i\le k}Z_i)^2}{k-1}.
\]

This recovers {[}EQ{]}, which proves the equal-mass subcase of {[}H19,
Conjecture 1.16{]} and the full four-cell dimension-three question
{[}H14, Conjecture 3{]}. At unequal masses, \(L_p\) is generally not
scalar on label space; the nonzero initial slope (15) supplies new
metric information that the uniform statement does not contain. The
universal bound (4) cannot follow merely by calibrating its quadratic
form on the model moments; (11)--(18) prove it for all feasible
functions.

The ordinary unweighted prescribed-mass squared-moment conjecture is
false {[}EX{]}. The present theorem uses the explicitly defined label
metric \(H_p^+\) and makes no contrary Euclidean assertion. It is a
first-moment functional inequality, not the full simplex-valued Gaussian
Bobkov or positive-correlation noise-stability statement. The profile
Hessian identity is prior work; its application here bounds all
functions, beyond variations of perimeter minimizers.

\section{6. Proof and verification
dependencies}\label{proof-and-verification-dependencies}

The non-elementary geometric input is {[}MN22, Theorem 1.1{]}, used only
with \(k=n+1\), its positive-mass hypotheses and its interface-perimeter
normalization. The established profile Hessian identity is stated with
its tangent-space convention and independently normalized in Section 2.
The remaining tools are Gaussian integration by parts, dominated
convergence, the implicit function theorem, matrix square roots,
pseudoinverses and one-variable calculus. There is no numerical or
solver premise, no differentiation at a singular endpoint, no assumption
of covariance concavity and no rank classification.

The package includes internal model reviews, finite exact diagnostics
and partial Lean proofs with explicit scope. They do not certify
Gaussian measure, the imported geometric theorem or the entire analytic
endpoint. Internal model reviews are not external human peer review. The
earlier equal-mass result is preserved as its own immutable record; this
theorem is a generalization of the same proof mechanism, not an
independent counting of the equal-mass case.


\section*{Source edition and reproduction}
This signed article edition retains the mathematical argument of the
previously archived source: \url{https://doi.org/10.5281/zenodo.23248386}.
The original source, exact inputs, checking programs and prior review
records remain available there. The companion README identifies the
frozen Git version and the scope of each auxiliary computation.
No computational certificate is needed for the written reduction.

\section*{Funding and research methods}
This research was conducted independently and received no external
funding. OpenAI AI assistants assisted mathematical research, verification
code and internal review. Mathematical claims are supported by the
written proof and the specified public evidence; AI output is not itself
proof evidence. Internal AI-assisted review is not external human peer
review.
\section*{Declaration of generative AI and AI-assisted technologies
in the writing process}
OpenAI AI assistants assisted drafting, organization and LaTeX editing.
The author takes responsibility for the manuscript and its descriptions
of the proof and verification evidence.

\section{References}\label{references}

\begin{itemize}
\tightlist
\item
  {[}MN22{]} E. Milman and J. Neeman, \emph{The Gaussian Double-Bubble
  and Multi-Bubble Conjectures}, Annals of Mathematics 195 (2022), no.
  1, 89--206; DOI 10.4007/annals.2022.195.1.2; arXiv:1805.10961v3,
  Theorem 1.1 and Proposition 2.4.
\item
  {[}H14{]} S. Heilman, \emph{Euclidean Partitions Optimizing Noise
  Stability}, Electronic Journal of Probability 19 (2014), no. 71,
  1--37; arXiv:1211.7138v2, Conjecture 3.
\item
  {[}H19{]} S. Heilman, \emph{Stable Gaussian Minimal Bubbles},
  arXiv:1901.03934v1, Problem 1.15 and Conjecture 1.16.
\item
  {[}EQ{]} mxym/math research project, \emph{Equal Gaussian cells: the
  sharp simplex first-moment theorem}, immutable release
  \nolinkurl{gaussian-balanced-simplex-all-k-v1}, commit
  \nolinkurl{ebfd90558439afca21babeb7b77c8d8837ea745f}, 2026.
\item
  {[}EX{]} mxym/math research project, \emph{A counterexample to the
  prescribed-mass Gaussian regular-simplex conjecture}, repository
  companion \nolinkurl{research/gaussian-fixed-mass-propeller-counterexample}, 2026.
\item
  {[}LS{]} The source/version and equivalence screen in this package's
  LITERATURE\_STATUS.md; limited source check of 8 October 2026, not a
  priority assessment.
\end{itemize}

\end{document}
